\documentclass{article}
\usepackage{amsmath,amssymb}
\usepackage{xurl}
\usepackage[hidelinks]{hyperref}
% Shared commands for notes in docs/paper-gaps/.

\DeclareUnicodeCharacter{2080}{\ifmmode{}_{0}\else\textsubscript{0}\fi}
\DeclareUnicodeCharacter{2081}{\ifmmode{}_{1}\else\textsubscript{1}\fi}
\DeclareUnicodeCharacter{2082}{\ifmmode{}_{2}\else\textsubscript{2}\fi}
\DeclareUnicodeCharacter{2099}{\ifmmode{}_{n}\else\textsubscript{n}\fi}

\newcommand{\Lean}{\textsc{Lean}}
\ExplSyntaxOn
\NewDocumentCommand{\leanid}{m}{
  \ifmmode
    \textsf{\detokenize{#1}}
  \else
    \group_begin:
    \urlstyle{sf}
    \tl_set:Nn \l_tmpa_tl {#1}
    \tl_replace_all:Nnn \l_tmpa_tl {\_} {_}
    \exp_args:NV \path \l_tmpa_tl
    \group_end:
  \fi
}
\ExplSyntaxOff
\providecommand{\tr}{\operatorname{tr}}
\newcommand{\tnleanIssueBase}{https://github.com/LionSR/TNLean/issues/}
\newcommand{\tnleanPrBase}{https://github.com/LionSR/TNLean/pull/}
\newcommand{\tnleanDocBase}{https://sirui-lu.com/TNLean/docs/}
\newcommand{\ghissue}[1]{\href{\tnleanIssueBase#1}{GitHub issue \##1}}
\newcommand{\ghpr}[1]{\href{\tnleanPrBase#1}{PR \##1}}
\newcommand{\doclink}[1]{\href{\tnleanDocBase#1.html}{\path{#1}}}

% Dirac notation and the identity operator, as in blueprint/src/macros/common.tex.
% \providecommand keeps note-local definitions authoritative.
\usepackage{dsfont}
\providecommand{\ket}[1]{|#1\rangle}
\providecommand{\bra}[1]{\langle #1|}
\providecommand{\braket}[2]{\langle #1 | #2 \rangle}
\providecommand{\ketbra}[2]{|#1\rangle\!\langle #2|}
\providecommand{\Id}{\mathds{1}}
\providecommand{\one}{\mathds{1}}
\providecommand{\C}{\mathbb{C}}
\providecommand{\MN}[1]{M_{#1}(\C)}
\providecommand{\MD}{M_D(\C)}

% External referencing for paper sources.  The build helper writes sanitized
% auxiliary files under build/paper-aux/ and excludes duplicated source labels.
\usepackage{xr}

% Import a generated paper auxiliary file with a caller-supplied label prefix.
% The candidate paths support builds from docs/paper-gaps/, the repository root,
% and build/paper-aux/ itself.
\newcommand{\importpaperaux}[2]{%
  \IfFileExists{../../build/paper-aux/#2.aux}{%
    \externaldocument[#1]{../../build/paper-aux/#2}%
  }{%
    \IfFileExists{build/paper-aux/#2.aux}{%
      \externaldocument[#1]{build/paper-aux/#2}%
    }{%
      \IfFileExists{#2.aux}{%
        \externaldocument[#1]{#2}%
      }{}%
    }%
  }%
}

% General external reference macros
\newcommand{\paperref}[2]{\ref{#1-#2}}
\newcommand{\papereqref}[2]{(\ref{#1-#2})}

% CPSV16 source (1606.00608 / MPDO-22-12-17-2.tex) external document and shortcuts
\importpaperaux{cpsv16-}{1606.00608/MPDO-22-12-17-2-tnlean}
\newcommand{\cpsvref}[1]{\paperref{cpsv16}{#1}}
\newcommand{\cpsveqref}[1]{\papereqref{cpsv16}{#1}}

% Verdict marker: \gapnote{<kind>}{<status>}. Kind is the policy
% classification (clarification, local-correction, scope-restriction,
% unfaithful, false-source, open-gap); status is open, wip (elimination
% underway), resolved, or historical. Machine-read by the site index;
% prints a verdict line.
\newcommand{\gapnote}[2]{\par\noindent\textbf{Verdict:} #1 (#2).\par}

\title{Quantum-Double PEPS: Normalization, Representations and Scope}
\author{}
\date{2026-09-26}
\begin{document}
\maketitle
\gapnote{scope-restriction}{open}

\section{What the sources print}
Definition~``def:2d-Ug-inj'' of \cite{Schuch2010PEPS} calls a PEPS tensor
$A$ $G$-injective for a semi-regular representation $g\mapsto U_g$, one
containing every irreducible representation of the finite group $G$, if
$\mathcal P(A)U_g=\mathcal P(A)$ and $\mathcal P(A)$ has a left inverse
$L$ with $L\,\mathcal P(A)=\Pi_U$, the projector onto the invariant
subspace.\footnote{Source traceability:
    \path{Papers/1001.3807/paper_v3.tex}, lines 1010--1013 and
    1278--1296.}
Definition~``def:iso:isopeps'' calls it $G$-isometric if moreover $U_g$ is
the left-regular representation $L_g\ket h=\ket{gh}$ and $\mathcal P(A)$,
restricted to its domain and range, is unitary.\footnote{Source
    traceability: \path{Papers/1001.3807/paper_v3.tex}, lines
    1692--1700.}
The quantum-double tensor
\begin{align}
    K=\sum_{p,q,r,s\in G}\ket{pq^{-1},qr^{-1},rs^{-1},sp^{-1}}\otimes
    \ket{p,q}\bra{r,s}
    \label{eq:rmp_peps_qd_tensor}
\end{align}
is then called $G$-isometric.\footnote{Source traceability:
    \path{Papers/1001.3807/paper_v3.tex}, lines 2904--2911.}
Appendix~A of \cite{Cirac2021Matrix}, under ``The Toric Code and quantum
double models'', gives the same tensor with plaquette colors on the
virtual legs and the differences of adjacent colors as physical spins,
states that it is $G$-isometric for the regular representation, and
states that the virtual labels of the primal tensor fuse to the identity,
so that the primal tensor has a symmetry under every irreducible
representation.\footnote{Source traceability:
    \path{Papers/2011.12127/TN-Review-main.tex}, lines 2450--2465.}

\section{Normalization}
Neither tensor is normalized. Each physical configuration of differences
arises from exactly $|G|$ colorings, those related by one global shift
$h_i\mapsto h_ig$. For vectors $x,y$ invariant under the shift,
\begin{align}
    \langle Kx,Ky\rangle=|G|\,\langle x,y\rangle .
    \label{eq:rmp_peps_qd_isometry}
\end{align}
So $K$ restricted to the invariant subspace is $\sqrt{|G|}$ times a
unitary onto its range, and $K/\sqrt{|G|}$ satisfies the printed
definition. The formal notion of $G$-isometry asks for
$\langle Kx,Ky\rangle=c\,\langle x,y\rangle$ with one constant $c>0$,
which the rescaling by $c^{-1/2}$ turns into the printed notion; the
quantum-double tensor satisfies it with $c=|G|$.

\section{Representations}
The formal notions of $G$-injectivity and $G$-isometry take the
representation as a parameter and do not require it to be semi-regular
or left-regular; each instance names the representation it uses. The
quantum-double instance shifts every color by right multiplication,
$h\mapsto hg$, which leaves the differences $hh'^{-1}$ unchanged also for
non-abelian $G$. On each leg this is the right-regular representation
$R_g\ket h=\ket{hg^{-1}}$, and the unitary $\ket h\mapsto\ket{h^{-1}}$ intertwines it with the left-regular
one. Both are semi-regular.

\section{Scope: statements about the whole network}
Two statements of the review concern the contracted lattice rather than
one tensor.
\begin{enumerate}
    \item The primal tensor's virtual symmetry under every irreducible
    representation $\pi$ is formalized as the single-tensor identity
    $A_{t,r,b,l}\,\pi(t)\pi(r)\pi(b)^{-1}\pi(l)^{-1}=A_{t,r,b,l}\,\mathbb 1$
    for every entry, which follows from $trb^{-1}l^{-1}=1$ on the support.
    The resulting matrix product operator symmetry of the contracted
    network is not formalized.
    \item The review states that the equal-weight superposition of
    plaquette colorings is the equal-weight superposition of the
    configurations obeying Gauss' law. On a plane, or on a disc with open
    boundary, every Gauss-law configuration is a difference of colorings.
    On a torus the differences of plaquette colorings are exactly the
    Gauss-law configurations of trivial holonomy around the two
    non-contractible cycles, so the statement holds for that sector only.
    It is not formalized. The review also prints the Gauss law as
    $g_1g_2g_3^{-1}g_4^{-1}=0$, where the identity element $1$ is
    meant.\footnote{Source traceability:
        \path{Papers/2011.12127/TN-Review-main.tex}, line 2452.}
\end{enumerate}

\section{Verdict and elimination plan}
The note is a scope restriction: the review's statements about the
contracted network are formalized only for one tensor or not at all. The
normalization and the representation are local corrections within that
restriction.
The normalization and the choice of regular representation are
conventions and need no further work. For the scope restrictions, place
the quantum-double tensors on the torus, prove the pulling-through
identity of the representation operators across the contracted network,
and state the coloring superposition as the trivial-holonomy sector of
the Gauss-law configurations.

\bibliographystyle{alpha}
\bibliography{references}
\end{document}
